% generated by R/50 (v2026a); csv: t3_slope_bootstrap.csv
\begin{table}[ht]
\caption{Separate versus joint calls: bootstrap comparison of the slopes}
\label{tab:indep:boot}
\footnotesize
\centering
\resizebox{\ifdim\width>\linewidth\linewidth\else\width\fi}{!}{%
\begin{tabular}{llllllrr}
\toprule
Model & Specification & Joint call & Separate calls & Separate minus joint & Separate / joint & N & Bootstrap draws \\
\midrule
GPT-4.1-mini & OLS & -0.239 [-0.351, -0.123] & -0.340 [-0.446, -0.233] & -0.100 [-0.198, -0.012] & 1.419 [1.041, 2.354] & 180 & 2,000 \\
GPT-4.1-mini & country + decade FE & -0.200 [-0.320, -0.058] & -0.315 [-0.426, -0.178] & -0.116 [-0.244, 0.013] & 1.578 [0.928, 4.331] & 180 & 2,000 \\
Sonnet 4.6 & OLS & -0.557 [-0.695, -0.378] & -0.614 [-0.736, -0.405] & -0.057 [-0.124, 0.039] & 1.103 [0.929, 1.263] & 47 & 2,000 \\
Sonnet 4.6 & country + decade FE & -0.411 [-1.009, 0.202] & -0.360 [-1.639, 0.765] & 0.050 [-1.200, 1.236] & 0.878 [-8.744, 9.483] & 47 & 1,879 \\
\bottomrule
\end{tabular}}
\par\smallskip
\begin{minipage}{\linewidth}\scriptsize Slope of liberalism (overall) on populism (overall), 0-10 scales; point estimate and 95\% percentile interval over manifesto bootstrap draws (draws with a degenerate fixed-effects fit are dropped). A ratio above 1 means the slope is steeper with separate calls, i.e. the joint call does not inflate the association.\end{minipage}
\end{table}
